\section*{Chapitre \ref{ch:g10:lewis-and-shape} --- Schémas de Lewis et géométrie des molécules}

\begin{solution}{exo:g10:lewis-and-shape:1}
Une paire d'\omterm{def:g9:inside-the-atom:nucleus}{électrons} mise en commun par deux \omterm{def:g7:atoms-and-molecules:atom}{atomes}. Une paire
d'\omterm{def:g10:electron-shells:valence}{électrons de valence} appartenant à un seul \omterm{def:g7:atoms-and-molecules:atom}{atome}, non partagée.
\end{solution}

\begin{solution}{exo:g10:lewis-and-shape:2}
H~: 1 (il lui manque un \omterm{def:g9:inside-the-atom:nucleus}{électron} pour son duet). C~: 4, N~: 3, O~: 2 (il
leur manque 4, 3 et 2 \omterm{def:g9:inside-the-atom:nucleus}{électrons} pour un octet).
\end{solution}

\begin{solution}{exo:g10:lewis-and-shape:3}
\chemfig{H-\charge{0=\:,90=\:,270=\:}{Cl}}~: le chlore possède le
\omterm{def:g10:lewis-and-shape:lone-pair}{doublet liant} et trois \omterm{def:g10:lewis-and-shape:lone-pair}{doublets non liants}, soit 8 \omterm{def:g9:inside-the-atom:nucleus}{électrons}.
\end{solution}

\begin{solution}{exo:g10:lewis-and-shape:4}
$5 + 3 \times 1 = 8$ \omterm{def:g9:inside-the-atom:nucleus}{électrons}, 4 doublets (trois liants, un non liant).
\end{solution}

\begin{solution}{exo:g10:lewis-and-shape:5}
\ce{CH4}~: tétraèdre. \ce{NH3}~: pyramide. \ce{H2O}~: coudée. \ce{CO2}~:
linéaire.
\end{solution}

\begin{solution}{exo:g10:lewis-and-shape:6}
\chemfig{H-\charge{90=\:,270=\:}{S}-H}~: comme l'eau, quatre groupes
autour du soufre, dont deux \omterm{def:g10:lewis-and-shape:lone-pair}{doublets non liants}~: coudée.
\end{solution}

\begin{solution}{exo:g10:lewis-and-shape:7}
$2 \times 5 = 10$ \omterm{def:g9:inside-the-atom:nucleus}{électrons}, 5 doublets. Une liaison simple et 4 \omterm{def:g10:lewis-and-shape:lone-pair}{doublets
non liants} laissent chaque azote avec 6 \omterm{def:g9:inside-the-atom:nucleus}{électrons}~; en transformant deux
\omterm{def:g10:lewis-and-shape:lone-pair}{doublets non liants} en \omterm{def:g10:lewis-and-shape:lone-pair}{doublets liants}, on obtient une \omterm{def:g10:lewis-and-shape:multiple-bond}{liaison triple} et
un \omterm{def:g10:lewis-and-shape:lone-pair}{doublet non liant} sur chaque \omterm{def:g7:atoms-and-molecules:atom}{atome}~:
\chemfig{\charge{180=\:}{N}~\charge{0=\:}{N}}.
\end{solution}

\begin{solution}{exo:g10:lewis-and-shape:8}
\chemfig{H-C(-[2]H)(-[6]H)-\charge{90=\:,270=\:}{O}-H}~: quatre liaisons
autour du carbone, tétraédrique~; autour de l'oxygène, deux liaisons et
deux \omterm{def:g10:lewis-and-shape:lone-pair}{doublets non liants}, coudée.
\end{solution}

\begin{solution}{exo:g10:lewis-and-shape:9}
Dans \ce{CO2}, le carbone a deux groupes (deux \omterm{def:g10:lewis-and-shape:multiple-bond}{liaisons doubles}) et aucun
\omterm{def:g10:lewis-and-shape:lone-pair}{doublet non liant}~: ils pointent dans des directions opposées. Dans
\ce{H2O}, l'oxygène a quatre groupes (deux liaisons, deux \omterm{def:g10:lewis-and-shape:lone-pair}{doublets non
liants}) en disposition tétraédrique~: les deux liaisons forment un angle.
\end{solution}

\begin{solution}{exo:g10:lewis-and-shape:10}
\chemfig{C(-[3]H)(-[5]H)=C(-[1]H)-[7]H}~: autour de chaque carbone,
trois groupes (une \omterm{def:g10:lewis-and-shape:multiple-bond}{liaison double}, deux liaisons simples)~: un triangle
plan, angles proches de \ang{120}.
\end{solution}

\begin{solution}{exo:g10:lewis-and-shape:11}
Le carbone au centre, quatre liaisons simples vers des \omterm{def:g7:atoms-and-molecules:atom}{atomes} de chlore
portant chacun trois \omterm{def:g10:lewis-and-shape:lone-pair}{doublets non liants}~: quatre groupes autour du
carbone, un tétraèdre.
\end{solution}

\begin{solution}{exo:g10:lewis-and-shape:12}
L'éthanol, \chemfig{H-C(-[2]H)(-[6]H)-C(-[2]H)(-[6]H)-\charge{90=\:,270=\:}{O}-H}, et
le méthoxyméthane,
\chemfig{H-C(-[2]H)(-[6]H)-\charge{90=\:,270=\:}{O}-C(-[2]H)(-[6]H)-H}.
\end{solution}

\begin{solution}{exo:g10:lewis-and-shape:13}
Un \omterm{def:g10:lewis-and-shape:lone-pair}{doublet non liant} prend un peu plus de place qu'un \omterm{def:g10:lewis-and-shape:lone-pair}{doublet liant} et
rapproche les liaisons. L'ammoniac a un \omterm{def:g10:lewis-and-shape:lone-pair}{doublet non liant}, l'eau deux~:
l'angle diminue du méthane à l'ammoniac puis à l'eau.
\end{solution}

\begin{solution}{exo:g10:lewis-and-shape:14}
\chemfig{H-C~\charge{0=\:}{N}}~: deux groupes autour du carbone,
linéaire. \chemfig{H-C(-[6]H)=\charge{45=\:,315=\:}{O}}~: trois groupes
autour du carbone, un triangle plan.
\end{solution}

\begin{solution}{exo:g10:lewis-and-shape:15}
$3 \times 6 = 18$ \omterm{def:g9:inside-the-atom:nucleus}{électrons}, 9 doublets~:
\chemfig{\charge{180=\:,270=\:}{O}=\charge{90=\:}{O}-\charge{0=\:,90=\:,270=\:}{O}}.
L'oxygène central a trois groupes (une \omterm{def:g10:lewis-and-shape:multiple-bond}{liaison double}, une liaison
simple, un \omterm{def:g10:lewis-and-shape:lone-pair}{doublet non liant})~: la \omterm{def:g7:atoms-and-molecules:molecule}{molécule} est coudée.
\end{solution}

\begin{solution}{pb:g10:lewis-and-shape:1}
\textbf{1.} \chemfig{H-C(-[2]H)(-[6]H)-H}, \chemfig{H-\charge{90=\:}{N}(-[6]H)-H},
\chemfig{H-\charge{90=\:,270=\:}{O}-H}.

\textbf{2.} Quatre groupes dans chacune.

\textbf{3.} Aucun pour le méthane, un pour l'ammoniac, deux pour l'eau.

\textbf{4.} Le carbone forme quatre liaisons~; l'oxygène en forme deux
(ses deux autres groupes sont des \omterm{def:g10:lewis-and-shape:lone-pair}{doublets non liants}, qu'une boîte de
modèles ne représente pas par des bâtonnets).

\textbf{5.} \chemfig{H-C(-[2]H)(-[6]H)-C(-[2]H)(-[6]H)-H}~: un bâtonnet
relie les deux boules de carbone, et il faut six boules d'hydrogène.

\textbf{6.} Tétraèdre, pyramide, coudée.

\textbf{7.} Les \omterm{def:g10:lewis-and-shape:lone-pair}{doublets non liants} prennent plus de place que les
\omterm{def:g10:lewis-and-shape:lone-pair}{doublets liants} et resserrent les liaisons~; l'eau, avec deux \omterm{def:g10:lewis-and-shape:lone-pair}{doublets
non liants}, a le plus petit angle.

\textbf{8.} Les trous d'une boule de carbone forment des angles de
\ang{109.5}~; le vrai angle de l'eau est \ang{104.5}~: une boule
d'oxygène doit avoir ses deux trous percés à \ang{104.5}.

\textbf{9.} \ang{180}~: deux groupes autour du carbone pointent dans des
directions opposées, et la \omterm{def:g7:atoms-and-molecules:molecule}{molécule} est linéaire.

\textbf{10.} Deux sommets qui ne sont pas sur une même arête sont des
sommets opposés d'une même face~: leur distance est la diagonale de la
face, $\sqrt{a^2 + a^2} = a\sqrt{2}$.

\textbf{11.} Le centre est au milieu d'une grande diagonale~:
$a\sqrt{3}/2$.

\textbf{12.} H--H~: $\sqrt{2} \approx \qty{1.41}{cm}$~; C--H~:
$\sqrt{3}/2 \approx \qty{0.87}{cm}$.

\textbf{13.} Chaque hydrogène est sur un sommet du cube et le carbone en
son centre~: tous les sommets sont à la même distance du centre.

\textbf{14.} Le triangle C\,H$_1$\,H$_2$ est isocèle (C\,H$_1$ =
C\,H$_2$)~; la droite qui joint son sommet principal au milieu de la
base est perpendiculaire à la base.

\textbf{15.} $\sin \widehat{\text{M\,C\,H}_1} = \dfrac{\text{M\,H}_1}{\text{C\,H}_1}
= \dfrac{a\sqrt{2}/2}{a\sqrt{3}/2} = \sqrt{2/3} \approx 0.816$.

\textbf{16.} Demi-angle $\approx \ang{54.7}$~; angle H--C--H $\approx
2 \times 54.7 = \ang{109.5}$.

\textbf{17.} C'est l'angle de \ang{109.5} du tétraèdre donné dans le
tableau.

\textbf{18.} Les trous doivent être percés à \ang{109.5} les uns des
autres.
\end{solution}
